% Generated draft: requires mathematical and source review.
\documentclass[11pt]{article}
\usepackage[margin=25mm]{geometry}
\usepackage{fontspec}
\setmainfont{Latin Modern Roman}
\usepackage{amsmath,amssymb,mathtools,unicode-math}
\setmathfont{Latin Modern Math}
\usepackage{enumitem,needspace,xurl,hyperref}
\setlength{\parindent}{0pt}
\setlength{\parskip}{5pt}
\setlength{\emergencystretch}{3em}
\newcommand{\N}{\mathbb{N}}
\newcommand{\Z}{\mathbb{Z}}
\newcommand{\Q}{\mathbb{Q}}
\newcommand{\R}{\mathbb{R}}
\newcommand{\C}{\mathbb{C}}
\newcommand{\F}{\mathbb{F}}
\newcommand{\A}{\mathbb{A}}
\newcommand{\PP}{\mathbb{P}}
\newcommand{\E}{\mathbb{E}}
\newcommand{\eps}{\varepsilon}
\newcommand{\dd}{\,\mathrm{d}}
\begin{document}
\section*{UnsolvedMath ID: 20002724}
\textbf{Source problem code:} catalogue entry
\subsection{Definitions and mathematical statement}
The exact canonical field stops after ``Problem 6. 3. (S. Fomin) Conjecture: The Fomin-Reading generalization of the associahedron \(\Delta_W^{(k)}\) (see''. It is not a complete conjecture. The supplied verification says one printed Problem 6.3 was split across canonical records 165--167. The aid's restoration reports topology, Cohen--Macaulay connectivity, and a polytopal-manifold clause; the complete restored wording is not actually included here. This draft preserves that source gap rather than attributing a fully invented conjecture to Fomin.

For a finite Coxeter group \(W\) of rank \(n\) and integer \(k\ge1\), let \(\Delta_W^{(k)}\) be the Fomin--Reading generalized cluster complex: its vertices are \(k\) colored copies of the positive roots and one copy of each negative simple root, with the standard generalized compatibility relation. The aid does not reproduce that relation. It reports that Fomin--Reading's Euler-characteristic formula and Athanasiadis--Tzanaki deletion-shellability imply that this \((n-1)\)-dimensional complex is \((k+1)\)-Cohen--Macaulay and
\[
 \Delta_W^{(k)}\simeq\bigvee^{\operatorname{Cat}^{(k-1)}(W)}S^{n-1}.
\]
For irreducible \(W\) with degrees \(d_i\) and Coxeter number \(h\), \(\operatorname{Cat}^{(m)}(W)=\prod_i(mh+d_i)/d_i\), with products over components in the reducible case. The source's printed \(k\)-CM request is thus reported answered in stronger form. The exact polytopal-manifold assertion is missing, and its general status was not verified as solved by the aid.

Fix a coefficient field. A pure complex is \(r\)-CM in the deletion sense if deleting any set of fewer than \(r\) vertices leaves a Cohen--Macaulay complex of the same dimension. The explicit Boolean family in the aid is
\[
 \Delta_{A_1^n}^{(k)}=D_{k+1}*\cdots*D_{k+1}\quad(n\text{ factors}),
\]
where each \(D_{k+1}\) is a discrete set of \(k+1\) vertices. This complex is exactly \((k+1)\)-CM, over every field: deleting at most \(k\) vertices leaves all factors nonempty and preserves the join dimension, while deleting one complete factor of \(k+1\) vertices lowers the dimension. Its wedge multiplicity is \(k^n\).

For \(n\ge2\) and \(k>1\), the aid's obstruction says that this complex cannot be the full \((n-1)\)-skeleton of any simplicial homology manifold, including a manifold with boundary. Its graph is complete \(n\)-partite, so no simplex in a putative ambient simplicial complex can have more than \(n\) vertices. The ambient dimension must therefore be \(n-1\), and it must equal the given skeleton. But each codimension-one face lies in \(k+1>2\) facets, whereas the homology-manifold link condition permits only one or two. The obstruction assumes equality with the full skeleton, not merely an embedding as a subcomplex.

Consequently any polytopal-cell manifold realization with that full skeleton would need nonsimplicial cells. This is a necessary restriction, not a construction or a general impossibility for all polytopal manifolds. It also does not settle the missing general source clause.
\subsection{Short English statement}
The canonical conjecture is truncated. The aid reports stronger known topology and deletion-CM results for generalized cluster complexes, and shows that the Boolean family cannot be the full skeleton of a simplicial homology manifold when \(n\ge2,k>1\). A nonsimplicial realization is not ruled out or constructed.
\subsection{Sources}
\begin{itemize}
\item Local registry, permanent UnsolvedMath ID \(20002724\), exact public code \texttt{catalogue entry}; canonical location \texttt{aim-probability-notes.json}, \texttt{notes[165]}.
\item S. Fomin, Problem 6.3, AIM \emph{Braid groups, clusters and free probability}; supplied reference \url{https://aimath.org/WWN/braidgroups/braidgroups.pdf}. The verification field reports the split across records 165--167. This record itself ends after ``see''; no PDF was accessed here and no neighboring record is silently substituted.
\item Fomin--Reading Euler-characteristic formula and Athanasiadis--Tzanaki deletion-shellability theorem are named by the aid without full titles, dates, theorem numbers, or bibliographic entries.
\item Supplied research attempt 1: \texttt{partial\_result}, run and registry status \texttt{partially\_solved}, rigor label \texttt{proved}, low novelty confidence; classification \texttt{PARTIAL-PROGRESS}, created and updated 14 August 2026. The aid is machine generated, not peer reviewed, and requires independent verification before citation.
\end{itemize}
\end{document}
